\documentclass[10pt,a4paper]{article}

% ---- Paquetes ----
\usepackage[utf8]{inputenc}
\usepackage[T1]{fontenc}
\usepackage[spanish]{babel}
\usepackage{amsmath,amssymb,amsfonts}
\usepackage[margin=1.8cm]{geometry}
\usepackage{enumitem}
\usepackage{hyperref}
\usepackage{physics}

% Compactar espaciado vertical
\setlength{\parskip}{2pt}
\setlength{\parindent}{0pt}
\usepackage{titlesec}
\titlespacing*{\section}{0pt}{6pt}{3pt}
\titlespacing*{\subsection}{0pt}{4pt}{2pt}

\title{\textbf{Examen --- Sistema de Ecuaciones Diferenciales}}
\author{}
\date{}

\begin{document}
\maketitle

% ============================================================
\section{Sistema}
% ============================================================

\[
\boxed{
\begin{cases}
\displaystyle \dv{x}{t} = 2x + y \\[2pt]
\displaystyle \dv{y}{t} = -x + 4y
\end{cases}
}
\]

\noindent\textbf{Condiciones iniciales (para Laplace):}\qquad
\(x(0) = 1,\quad y(0) = 0\).

% ============================================================
\section{Método de Eliminación (7 pasos)}
% ============================================================

\subsection{Paso 1 --- Despejar una variable}

De la primera ecuación:
\[
y = \dv{x}{t} - 2x = x' - 2x.
\]

\subsection{Paso 2 --- Derivar la expresión anterior}
\[
y' = x'' - 2x'.
\]

\subsection{Paso 3 --- Sustituir en la otra ecuación}

La segunda ecuación es \(y' = -x + 4y\). Sustituimos:
\[
x'' - 2x' = -x + 4(x' - 2x).
\]

\subsection{Paso 4 --- Obtener una EDO}
\[
\begin{aligned}
x'' - 2x' &= -x + 4x' - 8x \\
x'' - 2x' - 4x' + x + 8x &= 0 \\
\Rightarrow\quad x'' - 6x' + 9x &= 0.
\end{aligned}
\]

\subsection{Paso 5 --- Resolver la EDO}

Ecuación característica:
\[
r^{2} - 6r + 9 = (r - 3)^{2} = 0 \;\Longrightarrow\; r = 3\ \text{(doble)}.
\]

Solución general:
\[
x(t) = C_{1} e^{3t} + C_{2} t e^{3t}.
\]

\subsection{Paso 6 --- Sustituir para hallar la variable faltante}

De \(y = x' - 2x\):
\[
\begin{aligned}
x' &= 3C_{1} e^{3t} + C_{2} e^{3t} + 3C_{2} t e^{3t},\\[2pt]
y  &= \bigl(3C_{1} e^{3t} + C_{2} e^{3t} + 3C_{2} t e^{3t}\bigr)
      - 2\bigl(C_{1} e^{3t} + C_{2} t e^{3t}\bigr) \\[2pt]
   &= (3C_{1} - 2C_{1}) e^{3t} + C_{2} e^{3t} + (3C_{2} - 2C_{2}) t e^{3t} \\[2pt]
   &= C_{1} e^{3t} + C_{2} e^{3t} + C_{2} t e^{3t} \\[2pt]
   &= (C_{1} + C_{2}) e^{3t} + C_{2} t e^{3t}.
\end{aligned}
\]

\subsection{Paso 7 --- Escribir la solución completa}
\[
\boxed{
\begin{aligned}
x(t) &= C_{1} e^{3t} + C_{2} t e^{3t},\\[2pt]
y(t) &= (C_{1} + C_{2}) e^{3t} + C_{2} t e^{3t}.
\end{aligned}
}
\]

\smallskip
\noindent\textit{Nota:} Las constantes \(C_{1}, C_{2}\) se determinan si hay condiciones iniciales; aquí se dejan así porque eliminación no las usa.

% ============================================================
\section{Método del Operador Diferencial (9 pasos)}
% ============================================================

\subsection{Paso 1 --- Reordenar las ED dejando todo a la izquierda}
\[
\begin{aligned}
\dv{x}{t} - 2x - y &= 0 \quad\Longrightarrow\quad (D - 2)x - y = 0,\\[2pt]
\dv{y}{t} + x - 4y &= 0 \quad\Longrightarrow\quad x + (D - 4)y = 0.
\end{aligned}
\]

\subsection{Paso 2 --- Sustituir derivadas por \(D\) (ya está hecho)}

\subsection{Paso 3 --- Formar y calcular el determinante operacional}

Matriz del sistema:
\[
\begin{pmatrix}
D - 2 & -1 \\
1     & D - 4
\end{pmatrix}.
\]

Determinante:
\[
\Delta = (D - 2)(D - 4) - (-1)(1) = D^{2} - 6D + 8 + 1 = D^{2} - 6D + 9.
\]

\subsection{Paso 4 --- Obtener la ecuación característica}
\[
D^{2} - 6D + 9 = 0 \;\Longrightarrow\; r^{2} - 6r + 9 = 0.
\]

\subsection{Paso 5 --- Resolver raíces (caso 2: real repetida)}
\[
r = 3 \text{ (doble)}.
\]

\subsection{Paso 6 --- Proponer solución general para una variable}
\[
x(t) = C_{1} e^{3t} + C_{2} t e^{3t}.
\]

\subsection{Paso 7 --- Sustituir en ecuación original}

Usamos la primera ecuación original: \(y = (D - 2)x\),
\[
y = (D - 2)(C_{1} e^{3t} + C_{2} t e^{3t}) = Dx - 2x.
\]

Calculamos \(Dx = x'\):
\[
Dx = x' = 3C_{1} e^{3t} + C_{2} e^{3t} + 3C_{2} t e^{3t}.
\]

Entonces:
\[
\begin{aligned}
y &= \bigl[3C_{1} e^{3t} + C_{2} e^{3t} + 3C_{2} t e^{3t}\bigr]
     - 2\bigl[C_{1} e^{3t} + C_{2} t e^{3t}\bigr] \\[2pt]
  &= (3C_{1} - 2C_{1}) e^{3t} + C_{2} e^{3t} + (3C_{2} - 2C_{2}) t e^{3t} \\[2pt]
  &= C_{1} e^{3t} + C_{2} e^{3t} + C_{2} t e^{3t}.
\end{aligned}
\]

\subsection{Paso 8 --- Relacionar constantes}

Ya están relacionadas; \(y\) depende de \(C_{1}, C_{2}\) igual que antes:
\[
y = (C_{1} + C_{2}) e^{3t} + C_{2} t e^{3t}.
\]

\subsection{Paso 9 --- Escribir solución final}
\[
\boxed{
\begin{aligned}
x(t) &= C_{1} e^{3t} + C_{2} t e^{3t},\\[2pt]
y(t) &= (C_{1} + C_{2}) e^{3t} + C_{2} t e^{3t}.
\end{aligned}
}
\]

\smallskip
\noindent\textit{(Idéntica a eliminación.)}

% ============================================================
\section{Método de Transformada de Laplace (7 pasos)}
% ============================================================

\subsection{Paso 1 --- Aplicar Laplace a cada EDL}
\[
\begin{aligned}
\mathcal{L}\{x'\} &= sX(s) - x(0) = sX - 1,\\
\mathcal{L}\{y'\} &= sY(s) - y(0) = sY,\\
\mathcal{L}\{2x + y\} &= 2X + Y,\\
\mathcal{L}\{-x + 4y\} &= -X + 4Y.
\end{aligned}
\]

Las ecuaciones transformadas:
\[
\begin{aligned}
sX - 1 &= 2X + Y \;\Longrightarrow\; (s - 2)X - Y = 1,\\
sY    &= -X + 4Y \;\Longrightarrow\; X + (s - 4)Y = 0.
\end{aligned}
\]

\subsection{Paso 2 --- Aplicar condiciones iniciales (ya incluidas)}

\subsection{Paso 3 --- Ordenar el sistema algebraico}
\[
\begin{cases}
(s - 2)X - Y = 1,\\
X + (s - 4)Y = 0.
\end{cases}
\]

\subsection{Paso 4 --- Resolver el sistema para \(X(s)\) y \(Y(s)\)}

De la segunda: \(X = -(s - 4)Y\). Sustituimos en la primera:
\[
\begin{aligned}
(s - 2)\bigl[-(s - 4)Y\bigr] - Y &= 1 \\
-(s - 2)(s - 4)Y - Y &= 1 \\
-\bigl[(s - 2)(s - 4) + 1\bigr]Y &= 1 \\
-\bigl[(s^{2} - 6s + 8) + 1\bigr]Y &= -\bigl[s^{2} - 6s + 9\bigr]Y = 1 \\
-(s - 3)^{2} Y &= 1.
\end{aligned}
\]

Por tanto,
\[
Y(s) = -\frac{1}{(s - 3)^{2}}.
\]

Luego:
\[
X(s) = -(s - 4)Y = -(s - 4)\left(-\frac{1}{(s - 3)^{2}}\right)
      = \frac{s - 4}{(s - 3)^{2}}.
\]

\subsection{Paso 5 --- Simplificar las ecuaciones}

Descomponemos \(X(s)\) en fracciones parciales:
\[
\frac{s - 4}{(s - 3)^{2}} = \frac{1}{s - 3} - \frac{1}{(s - 3)^{2}},
\]
(verificación: \((1)(s - 3) - 1 = s - 4\)).

\subsection{Paso 6 --- Aplicar transformada inversa de Laplace}
\[
\begin{aligned}
x(t) &= \mathcal{L}^{-1}\!\left\{\frac{1}{s - 3}\right\}
      - \mathcal{L}^{-1}\!\left\{\frac{1}{(s - 3)^{2}}\right\}
      = e^{3t} - t e^{3t},\\[3pt]
y(t) &= \mathcal{L}^{-1}\!\left\{-\frac{1}{(s - 3)^{2}}\right\}
      = -t e^{3t}.
\end{aligned}
\]

\subsection{Paso 7 --- Escribir solución final}
\[
\boxed{
\begin{aligned}
x(t) &= e^{3t} - t e^{3t},\\[2pt]
y(t) &= -t e^{3t}.
\end{aligned}
}
\]

\medskip
\noindent\textit{Nota:} Esta es la solución particular con las condiciones iniciales dadas. Comparando con la solución general de eliminación, si \(C_{1}=1\) y \(C_{2}=-1\) se obtiene exactamente este resultado.

\end{document}
